% GENERATED FILE. Edit canonical lessons, then run npm run generate:books.
% canonical-source-hash: fnv1a64:61b87926e77d45fa
% canonical-lessons: GE-C517-schauspieler, GE-C517-regisseur, GE-C517-publikum, GE-C517-klatschen, GE-C517-applaus

\chapter{Im Theater}
\label{ch:ge-im-theater}
\hlchaptermodality{517}

\begin{chapteropening}
\textbf{By the end of this chapter:} \emph{I can talk about plays, actors and audiences.}

Complete the last lesson of chapter 517: I can talk about plays, actors and audiences.
\end{chapteropening}

\section[der Schauspieler]{der Schauspieler — actor}

\label{lesson:GE-C517-schauspieler}



\begin{quote}
Before the new one: say the German for to perform, then the German for appearance.
\end{quote}

\subsection*{You'll want to know: der Schauspieler}
\textbf{der Schauspieler} — \textquotedblleft{}actor\textquotedblright{}.

\begin{quote}
Der Schauspieler hat die Rolle sehr gut gespielt.
\end{quote}

\begin{quote}
\emph{The actor played the role very well.}
\end{quote}

Plural: \emph{die Schauspieler}. A woman is \emph{die Schauspielerin}.

\begin{cousinweb}
\emph{Schauspieler} is built on \emph{Schauspiel} \textquotedblleft{}play\textquotedblright{}, from \emph{schauen} \textquotedblleft{}to look\textquotedblright{}, inherited from Old High German \emph{scouwōn}, and \emph{Spiel} \textquotedblleft{}game, performance\textquotedblright{} (EtymDB).
\end{cousinweb}

\begin{grammarlens}[title={What you've built}]
One more word for telling, explaining and describing at B1.
\end{grammarlens}

\subsection*{Your turn}
\begin{itemize}
\raggedright
  \item \emph{Say it:} \emph{der Schauspieler}
  \item \emph{Say it:} \emph{der Schauspieler}, once more
  \item \emph{Say it:} a sentence of your own with \emph{der Schauspieler}
  \item \emph{Recall:} say the German for to perform, then the German for appearance, then say \emph{der Schauspieler} again
\end{itemize}

\begin{tcolorbox}[breakable,colback=teal!4,colframe=teal!35!black,title={Before you move on}]
What does der Schauspieler mean? (\textquotedblleft{}Actor\textquotedblright{}.) Say it once more.
\end{tcolorbox}
\section[der Regisseur]{der Regisseur — director (of a film or play)}

\label{lesson:GE-C517-regisseur}



\begin{quote}
Before the new one: say the German for appearance, then the German for actor.
\end{quote}

\subsection*{You'll want to know: der Regisseur}
\textbf{der Regisseur} — \textquotedblleft{}director (of a film or play)\textquotedblright{}.

\begin{quote}
Der Regisseur arbeitet seit Jahren mit denselben Schauspielern.
\end{quote}

\begin{quote}
\emph{The director has worked with the same actors for years.}
\end{quote}

Plural: \emph{die Regisseure}. Said the French way: \emph{re-shi-SÖR}.

\begin{cousinweb}
\emph{Regisseur} was borrowed from French \emph{régisseur} \textquotedblleft{}manager\textquotedblright{} (EtymDB), built on \emph{régir} \textquotedblleft{}to manage\textquotedblright{}, from Latin \emph{regere} \textquotedblleft{}to rule\textquotedblright{}.
\end{cousinweb}

\begin{grammarlens}[title={What you've built}]
One more word for telling, explaining and describing at B1.
\end{grammarlens}

\subsection*{Your turn}
\begin{itemize}
\raggedright
  \item \emph{Say it:} \emph{der Regisseur}
  \item \emph{Say it:} \emph{der Regisseur}, once more
  \item \emph{Say it:} a sentence of your own with \emph{der Regisseur}
  \item \emph{Recall:} say the German for appearance, then the German for actor, then say \emph{der Regisseur} again
\end{itemize}

\begin{tcolorbox}[breakable,colback=teal!4,colframe=teal!35!black,title={Before you move on}]
What does der Regisseur mean? (\textquotedblleft{}Director (of a film or play)\textquotedblright{}.) Say it once more.
\end{tcolorbox}
\section[das Publikum]{das Publikum — audience}

\label{lesson:GE-C517-publikum}



\begin{quote}
Before the new one: say the German for actor, then the German for director.
\end{quote}

\subsection*{You'll want to know: das Publikum}
\textbf{das Publikum} — \textquotedblleft{}audience\textquotedblright{}.

\begin{quote}
Das Publikum lachte über jeden Witz.
\end{quote}

\begin{quote}
\emph{The audience laughed at every joke.}
\end{quote}

No plural. Singular in German: \emph{Das Publikum war begeistert.}

\begin{cousinweb}
\emph{Publikum} comes from Medieval Latin \emph{publicum} \textquotedblleft{}the public\textquotedblright{}, from Latin \emph{pūblicus} \textquotedblleft{}of the people\textquotedblright{} (Lewis \& Short).
\end{cousinweb}

\begin{grammarlens}[title={What you've built}]
One more word for telling, explaining and describing at B1.
\end{grammarlens}

\subsection*{Your turn}
\begin{itemize}
\raggedright
  \item \emph{Say it:} \emph{das Publikum}
  \item \emph{Say it:} \emph{das Publikum}, once more
  \item \emph{Say it:} a sentence of your own with \emph{das Publikum}
  \item \emph{Recall:} say the German for actor, then the German for director, then say \emph{das Publikum} again
\end{itemize}

\begin{tcolorbox}[breakable,colback=teal!4,colframe=teal!35!black,title={Before you move on}]
What does das Publikum mean? (\textquotedblleft{}Audience\textquotedblright{}.) Say it once more.
\end{tcolorbox}
\section[klatschen]{klatschen — to applaud; to gossip}

\label{lesson:GE-C517-klatschen}



\begin{quote}
Before the new one: say the German for director, then the German for audience.
\end{quote}

\subsection*{You'll want to know: klatschen}
\textbf{klatschen} — \textquotedblleft{}to applaud; to gossip\textquotedblright{}.

\begin{quote}
Am Ende klatschten alle begeistert.
\end{quote}

\begin{quote}
\emph{At the end everyone applauded enthusiastically.}
\end{quote}

\emph{Beifall klatschen} — to applaud. Informal: \emph{über die Nachbarn klatschen} — to gossip about the neighbours.

\begin{cousinweb}
\emph{klatschen} is usually explained as an imitative word for a slapping sound; EtymDB lists it, glossed \textquotedblleft{}to smack, to slap\textquotedblright{}, without a source.
\end{cousinweb}

\begin{grammarlens}[title={What you've built}]
One more word for telling, explaining and describing at B1.
\end{grammarlens}

\subsection*{Your turn}
\begin{itemize}
\raggedright
  \item \emph{Say it:} \emph{klatschen}
  \item \emph{Say it:} \emph{klatschen}, once more
  \item \emph{Say it:} a sentence of your own with \emph{klatschen}
  \item \emph{Recall:} say the German for director, then the German for audience, then say \emph{klatschen} again
\end{itemize}

\begin{tcolorbox}[breakable,colback=teal!4,colframe=teal!35!black,title={Before you move on}]
What does klatschen mean? (\textquotedblleft{}To applaud; to gossip\textquotedblright{}.) Say it once more.
\end{tcolorbox}
\section[der Applaus]{der Applaus — applause}

\label{lesson:GE-C517-applaus}



\begin{quote}
Before the new one: say the German for audience, then the German for to applaud.
\end{quote}

\subsection*{You'll want to know: der Applaus}
\textbf{der Applaus} — \textquotedblleft{}applause\textquotedblright{}.

\begin{quote}
Der Applaus dauerte fünf Minuten.
\end{quote}

\begin{quote}
\emph{The applause lasted five minutes.}
\end{quote}

No plural. Also \emph{der Beifall}.

\begin{cousinweb}
\emph{Applaus} comes from Latin \emph{applausus} \textquotedblleft{}applause\textquotedblright{} (EtymDB), from \emph{applaudere} \textquotedblleft{}to strike against\textquotedblright{}, built on \emph{plaudere} \textquotedblleft{}to strike\textquotedblright{}.
\end{cousinweb}

\begin{grammarlens}[title={What you've built}]
One more word for telling, explaining and describing at B1.
\end{grammarlens}

\subsection*{Your turn}
\begin{itemize}
\raggedright
  \item \emph{Say it:} \emph{der Applaus}
  \item \emph{Say it:} \emph{der Applaus}, once more
  \item \emph{Say it:} a sentence of your own with \emph{der Applaus}
  \item \emph{Recall:} say the German for audience, then the German for to applaud, then say \emph{der Applaus} again
\end{itemize}

\begin{tcolorbox}[breakable,colback=teal!4,colframe=teal!35!black,title={Before you move on}]
What does der Applaus mean? (\textquotedblleft{}Applause\textquotedblright{}.) Say it once more.
\end{tcolorbox}
